\providecommand{\pending}[1]{\textcolor{gray}{\textsc{#1}}}
\begin{abstract}
Benchmarks for automated experimental design make experiments instantaneous, free, and
infinitely repeatable. Real experiments are none of these. We introduce \slowlab, an
executable benchmark under four conditions that hold at once: feedback is slow, noise is
irreducible, trials are priced in the objective's own units, and commitments cannot be
rolled back. An agent runs a greenhouse --- it designs experiments, commits them, waits out
a growing cycle, and repeats. Ground truth is a published mechanistic crop model, so the
difficulty is not ours to tune, and everything is denominated in net profit per square metre
per day. Beside realised regret we report the regret a perfect reasoner would still carry
given the same design, computed from an instance distribution no agent sees. The gap
separates a campaign that could not have answered the question from one that could and did
not, and scores an agent against what was achievable rather than against a baseline we
happened to implement.

Across five models and four tasks, $99\%$ of designs on the six-factor screen cannot
separate the effects they were commissioned to separate, because most use two distinct
treatments replicated across the facility: under a hard parallelism cap, precision and
coverage compete for the same physical units. This is the opposite of the failure we
expected. Supplying standard tools then drives the two measures apart --- the designs
improve and the campaigns end worse in sixteen of twenty cells, because the inference
tool is fitted to the agent's own design, inherits its deficiency, and is deferred to in
proportion to how wrong it is. Under these conditions an agent cannot afford the cycle it
would take to find that out.
\end{abstract}
